\section{Introduction}
Vineyard yield estimates guide decisions on labour, logistics and crop thinning, yet they are still commonly made by destructive hand sampling, which is labour intensive and too sparse to capture variation across a block~\cite{nuske2014}. A camera carried along the rows could sample every vine. The most accessible version is a smartphone fixed to a vehicle that already drives the rows, such as a tractor or utility vehicle. It needs no depth sensor, LiDAR or multi-camera rig.

The obstacle is occlusion. Leaves, shoots and neighbouring bunches hide fruit, so even a perfect single-image detector undercounts. Hemming \etal~\cite{hemming2014} showed how much viewpoint matters: for sweet pepper, no single camera position detected more than 69\% of fruit, while five positions together reached 90\%. Occlusion is therefore not only a detector problem but a sampling problem. The question is which views to collect and how to combine them.

Harvesting robotics answers this with \emph{active perception}. A manipulator moves the camera to a new viewpoint for each fruit, by next-best-view planning or visual servoing, until the fruit is visible~\cite{magalhaes2022,burusa2024}. This suits a robot that must reach individual fruit, but it needs an arm, per-target planning and time at each plant. Vineyard monitoring is different. Vines are grown in straight trellised rows, and the vehicle's path is fixed by the row. The camera cannot choose its viewpoint. Instead, as the vehicle moves, it passively gathers a dense, ordered sweep of lateral views, in which the near foliage shifts more than the fruit behind it because of parallax. Existing vehicle-based vineyard systems use such sequences mainly to count visible fruit and then correct for hidden fruit statistically~\cite{nuske2014,millan2018}, to avoid double counting across frames~\cite{santos2020,liu2019}, or to predict hidden berries with learned models~\cite{quinones2025gcnet,kierdorf2022}. Tracking-based counters~\cite{liu2019,santos2020,stein2016} credit a fruit that is detected whole in some view, and FruitNeRF~\cite{meyer2024fruitnerf} counts fruit in a radiance field built from many posed views. As far as we know, none composes the partial observations of the sweep into a single de-occluded view of the fruit, so fruit never fully exposed in any one view stays uncounted.

We address that gap. Given a sequence from one moving RGB camera, our pipeline recovers grape pixels hidden in a reference view, using only colours observed in other views. Because it is purely image based, it needs no depth sensor or calibrated poses: alignment is estimated from fruit features in the images themselves, and the only motion input is the nominal view spacing used by a baseline rule that was inactive at the spacings tested (Sec.~\ref{sec:ablation}). As a controlled first step, we use a phone on a horizontal rail as a surrogate for vehicle motion.

Our contributions are (i) a fruit-only registration and observation-backed fusion pipeline: it segments fruit and foliage with SAM~3~\cite{carion2025}, aligns views on fruit features with a jointly solved and validated similarity model, and keeps a pixel only when several views spanning a minimum baseline agree on its colour; (ii) an evaluation that judges recovery by a downstream task as well as by pixels, applying an independently trained grape counter, GrapeSAM~\cite{tong2026vivid}, unchanged to the occluded view, the reconstruction and a leaf-free photograph; and (iii) a controlled study of viewpoint spacing (20--60~mm) and occlusion (3--6 foreground branches), including where and why recovery fails.
We also briefly explore separating red and green fruit, a step toward the ripe and unripe estimate that the title names. This report evaluates the occlusion-removal step only; it does not estimate yield, fruit mass, ripeness or vine health.
